\documentclass[11pt]{article}

\usepackage[margin=1in]{geometry}
\usepackage{amsmath, amssymb}
\usepackage{enumitem}
\usepackage{array}
\usepackage{fancyhdr}

\pagestyle{fancy}
\fancyhf{}
\lhead{Scientific Notation Practice (10.5--10.7)}
\rhead{Name:\ \rule{2.7in}{0.4pt}}
\cfoot{\thepage}
\setlength{\parindent}{0pt}

\newcommand{\work}[1]{\vspace{#1}\par}

\begin{document}

\begin{center}
{\Large \textbf{Practice Worksheet: Sections 10.5--10.7}}\\[4pt]
{\small 10.5 Reading Scientific Notation \quad 10.6 Writing Scientific Notation \quad 10.7 Operations in Scientific Notation}
\end{center}

\textbf{Directions:} Show your work in the space provided. For Section D, write each final answer in \textbf{scientific notation}.

\vspace{0.4em}\hrule\vspace{0.9em}

% =========================================================
% 10.5 Reading Scientific Notation
% =========================================================
\textbf{A. Tell whether the number is written in scientific notation. Explain. (Section 10.5)}

\begin{enumerate}[label=\arabic*.]
\item $18 \times 10^{4}$
\work{1.0in}

\item $0.93 \times 10^{-6}$
\work{1.0in}

\item $7.01 \times 10^{0}$
\work{1.0in}

\item $12.4 \times 10^{-3}$
\work{1.0in}
\end{enumerate}

\vspace{0.2em}\hrule\vspace{0.9em}

% =========================================================
% 10.5 Standard Form from Scientific Notation
% =========================================================
\textbf{B. Write the number in standard form. (Section 10.5)}

\begin{enumerate}[label=\arabic*., resume]
\item $3.6 \times 10^{7}$
\work{1.1in}

\item $5.08 \times 10^{-4}$
\work{1.1in}

\item $9.12 \times 10^{2}$
\work{1.1in}

\item $1.47 \times 10^{-1}$
\work{1.1in}
\end{enumerate}

\vspace{0.2em}\hrule\vspace{0.9em}

% =========================================================
% 10.6 Writing Scientific Notation
% =========================================================
\textbf{C. Write the number in scientific notation. (Section 10.6)}

\begin{enumerate}[label=\arabic*., resume]
\item $0.0000835$
\work{1.1in}

\item $46,\!300,\!000$
\work{1.1in}

\item $0.01209$
\work{1.1in}

\item $905,\!000$
\work{1.1in}
\end{enumerate}

\vspace{0.2em}\hrule\vspace{0.9em}

% =========================================================
% 10.7 Operations in Scientific Notation
% =========================================================
\textbf{D. Evaluate the expression. Write your answer in scientific notation. (Section 10.7)}

\begin{enumerate}[label=\arabic*., resume]
\item $(6.45\times 10^{5}) + (2.8\times 10^{5})$
\work{1.4in}

\item $(3.2\times 10^{-6}) - (7.5\times 10^{-7})$
\work{1.4in}

\item $(4.6\times 10^{3}) \times (1.9\times 10^{-2})$
\work{1.6in}

\item $(7.2\times 10^{-4}) \div (9.0\times 10^{-8})$
\work{1.6in}

\item $(9.1\times 10^{7}) + (3.7\times 10^{6})$
\work{1.6in}

\item $(5.04\times 10^{-3}) - (1.8\times 10^{-4})$
\work{1.6in}
\end{enumerate}

\vspace{0.2em}\hrule\vspace{0.9em}

% =========================================================
% Word problems + table (quiz-style)
% =========================================================
\textbf{E. Applications (Sections 10.5--10.7)}

\textbf{1. PLANETS.} The table shows approximate equatorial radii of several planets.

\vspace{0.4em}
\begin{tabular}{|>{\raggedright\arraybackslash}p{2.1in}|>{\raggedleft\arraybackslash}p{2.6in}|}
\hline
\textbf{Planet} & \textbf{Equatorial Radius (km)}\\
\hline
Mercury & $2.44\times 10^{3}$\\
Venus   & $6.05\times 10^{3}$\\
Earth   & $6.37\times 10^{3}$\\
Mars    & $3.39\times 10^{3}$\\
Jupiter & $7.15\times 10^{4}$\\
Saturn  & $6.03\times 10^{4}$\\
\hline
\end{tabular}

\begin{enumerate}[label=\alph*.]
\item Which planet has the \textbf{second-smallest} equatorial radius? Explain how you know.
\work{1.3in}

\item Which planet has the \textbf{second-largest} equatorial radius? Explain how you know.
\work{1.3in}
\end{enumerate}

\vspace{0.6em}

\textbf{2. SPACE DISTANCE.} A spacecraft is about $4\times 10^{5}$ kilometers from Earth.
If $1$ kilometer is $10^{3}$ meters, how far is the spacecraft from Earth in meters?
Write your answer in scientific notation.
\work{1.8in}

\textbf{3. SKIN THICKNESS.} A layer of skin is about $0.00042$ meters thick.
Write this number in scientific notation.
\work{1.4in}

\textbf{4. ORBITS.} The Sun takes about $2.25\times 10^{8}$ years to orbit the center of the Milky Way.
A dwarf planet takes about $3.0\times 10^{2}$ years to orbit the Sun.
How many times does the dwarf planet orbit the Sun while the Sun completes one orbit around the Milky Way?
Write your answer in standard form.
\work{2.0in}

\end{document}